\documentclass[12pt,a4paper]{article}

% ---- Encoding & Font ----
\usepackage{iftex}
\ifPDFTeX
  \usepackage[utf8]{inputenc}
  \usepackage[T1]{fontenc}
\else
  \usepackage{fontspec}  % XeTeX/LuaTeX (tectonic): Unicode nativo (·, —, ¿, ª, §)
\fi
\usepackage[spanish]{babel}
\usepackage{csquotes}

% ---- Page layout ----
\usepackage[top=2.5cm, bottom=2.5cm, left=2.5cm, right=2.5cm]{geometry}
\usepackage{setspace}
\onehalfspacing

% ---- Math ----
\usepackage{amsmath, amssymb, amsthm}
\usepackage{mathtools}
\usepackage{physics}

% ---- Graphics ----
\usepackage{graphicx}
\usepackage{tikz}
\usepackage{float}
\usepackage{caption}

% ---- Tables ----
\usepackage{array}
\usepackage{booktabs}
\usepackage{tabularx}
\usepackage{colortbl}

% ---- Colours ----
\usepackage{xcolor}
\definecolor{notebg}{HTML}{FFF3CD}
\definecolor{noteborder}{HTML}{FFC107}
\definecolor{boxred}{HTML}{F8D7DA}
\definecolor{boxredborder}{HTML}{F5C6CB}

% ---- Boxes ----
\usepackage{mdframed}
\usepackage{tcolorbox}
\tcbuselibrary{most}

\newtcolorbox{keybox}[1][]{
  colback=blue!5,
  colframe=blue!70!black,
  arc=4pt,
  boxrule=1pt,
  fonttitle=\bfseries,
  title={#1}
}

\newtcolorbox{warnbox}[1][]{
  colback=boxred,
  colframe=boxredborder,
  coltitle=black!75,
  arc=4pt,
  boxrule=1pt,
  fonttitle=\bfseries,
  title={#1}
}

\newtcolorbox{notebox}[1][]{
  colback=notebg,
  colframe=noteborder,
  coltitle=black!75,
  arc=4pt,
  boxrule=1pt,
  fonttitle=\bfseries,
  title={#1}
}

% ---- Diagramas (gráficas) ----
\usepackage{pgfplots}
\pgfplotsset{compat=1.18}
\usetikzlibrary{babel}  % babel español activa `>`; esto evita romper las flechas `->`
% courses/CDIV2017/Clases/assets/tikzstyles.tex
% ---------------------------------------------------------------------------
% Estilos compartidos para los diagramas del curso Cálculo 1 (CDIV2017).
% Fuente única del "look" visual (colores, grosores, escalas) para que todas
% las figuras (TikZ / pgfplots: conjuntos, rectas numéricas, gráficas de
% funciones) sean consistentes entre clases.
%
% Uso: la clase debe cargar antes `tikz` y `pgfplots` (bloque §6.3 del
%      CLAUDE.md del curso), y luego
%      % courses/CDIV2017/Clases/assets/tikzstyles.tex
% ---------------------------------------------------------------------------
% Estilos compartidos para los diagramas del curso Cálculo 1 (CDIV2017).
% Fuente única del "look" visual (colores, grosores, escalas) para que todas
% las figuras (TikZ / pgfplots: conjuntos, rectas numéricas, gráficas de
% funciones) sean consistentes entre clases.
%
% Uso: la clase debe cargar antes `tikz` y `pgfplots` (bloque §6.3 del
%      CLAUDE.md del curso), y luego
%      % courses/CDIV2017/Clases/assets/tikzstyles.tex
% ---------------------------------------------------------------------------
% Estilos compartidos para los diagramas del curso Cálculo 1 (CDIV2017).
% Fuente única del "look" visual (colores, grosores, escalas) para que todas
% las figuras (TikZ / pgfplots: conjuntos, rectas numéricas, gráficas de
% funciones) sean consistentes entre clases.
%
% Uso: la clase debe cargar antes `tikz` y `pgfplots` (bloque §6.3 del
%      CLAUDE.md del curso), y luego
%      \input{../assets/tikzstyles.tex}
% (compilar desde el directorio de la clase para que la ruta relativa resuelva).
%
% Paleta propia de este curso (no se reusa la de courses/Fisica3-2015): mismos
% tonos base (azul/rojo/ámbar) para alinear con keybox/notebox/warnbox, pero
% archivo independiente, sin dependencia cruzada entre cursos.
% ---------------------------------------------------------------------------

% ---- Paleta (alineada con las cajas del preámbulo: keybox/notebox/warnbox) ----
\definecolor{figblue}{HTML}{1F4E79}   % azul principal (líneas, keybox, conjuntos)
\definecolor{figred}{HTML}{B03A2E}    % rojo de énfasis (warnbox, región resaltada)
\definecolor{figamber}{HTML}{B8860B}  % ámbar (notebox)
\definecolor{figgray}{HTML}{555555}   % auxiliares, ejes, guías, universo

% ---- TikZ: librerías y estilos reutilizables ----
% Nota: las puntas se declaran como `-{Latex[...]}` y no como `->`. La forma
% con llaves NO contiene un `>` literal, así que es inmune al `>` activo que
% introduce babel español (de todos modos se carga \usetikzlibrary{babel} en
% el bloque §6.3 para cubrir cualquier `->` suelto).
\usetikzlibrary{arrows.meta,calc,angles,quotes,patterns}

\tikzset{
  figline/.style={figblue, line width=0.9pt},
  figaux/.style={figgray, dashed, line width=0.6pt},
  figaxis/.style={figgray, line width=0.7pt, -{Latex[length=2.2mm,width=1.5mm]}},
  % conjuntos (diagramas de Venn / patatoidales)
  figset/.style={figline, fill=figblue!6},
  figsetB/.style={figred, line width=0.9pt, fill=figred!6},
  figsetC/.style={figamber, line width=0.9pt, fill=figamber!10},
  figuniverso/.style={figgray, line width=0.7pt},
  % puntos de la recta / conjuntos numéricos
  figpto/.style={circle, inner sep=0pt, minimum size=5pt, draw=none},
  figptoabierto/.style={circle, inner sep=0pt, minimum size=5pt, draw=figblue, line width=0.9pt, fill=white},
  figptocerrado/.style={circle, inner sep=0pt, minimum size=5pt, draw=figblue, line width=0.9pt, fill=figblue},
  figptoY/.style={figpto, fill=figblue},
  figptoZ/.style={figpto, fill=figred},
  % segmentos gruesos para intervalos sobre la recta
  figintervalo/.style={figblue, line width=2.2pt},
  % vectores/flechas de transformación (traslación, dilatación)
  figvec/.style={figred, line width=1pt, -{Latex[length=2.2mm,width=1.6mm]}},
  % etiquetas
  figlbl/.style={font=\small},
  figlblaux/.style={font=\footnotesize, figgray},
}

% ---- pgfplots: estilo base de las gráficas de funciones ----
\pgfplotsset{
  calcfig/.style={
    width=10cm, height=6cm,
    axis lines=middle,
    xlabel style={font=\small, at={(axis description cs:1,0.5)}, anchor=west},
    ylabel style={font=\small, at={(axis description cs:0.5,1)}, anchor=south},
    tick label style={font=\footnotesize},
    every axis plot/.append style={line width=1.1pt, figblue},
    grid=none,
    legend cell align=left,
    legend style={font=\footnotesize, draw=figgray!40},
    clip=false,
  },
}

% (compilar desde el directorio de la clase para que la ruta relativa resuelva).
%
% Paleta propia de este curso (no se reusa la de courses/Fisica3-2015): mismos
% tonos base (azul/rojo/ámbar) para alinear con keybox/notebox/warnbox, pero
% archivo independiente, sin dependencia cruzada entre cursos.
% ---------------------------------------------------------------------------

% ---- Paleta (alineada con las cajas del preámbulo: keybox/notebox/warnbox) ----
\definecolor{figblue}{HTML}{1F4E79}   % azul principal (líneas, keybox, conjuntos)
\definecolor{figred}{HTML}{B03A2E}    % rojo de énfasis (warnbox, región resaltada)
\definecolor{figamber}{HTML}{B8860B}  % ámbar (notebox)
\definecolor{figgray}{HTML}{555555}   % auxiliares, ejes, guías, universo

% ---- TikZ: librerías y estilos reutilizables ----
% Nota: las puntas se declaran como `-{Latex[...]}` y no como `->`. La forma
% con llaves NO contiene un `>` literal, así que es inmune al `>` activo que
% introduce babel español (de todos modos se carga \usetikzlibrary{babel} en
% el bloque §6.3 para cubrir cualquier `->` suelto).
\usetikzlibrary{arrows.meta,calc,angles,quotes,patterns}

\tikzset{
  figline/.style={figblue, line width=0.9pt},
  figaux/.style={figgray, dashed, line width=0.6pt},
  figaxis/.style={figgray, line width=0.7pt, -{Latex[length=2.2mm,width=1.5mm]}},
  % conjuntos (diagramas de Venn / patatoidales)
  figset/.style={figline, fill=figblue!6},
  figsetB/.style={figred, line width=0.9pt, fill=figred!6},
  figsetC/.style={figamber, line width=0.9pt, fill=figamber!10},
  figuniverso/.style={figgray, line width=0.7pt},
  % puntos de la recta / conjuntos numéricos
  figpto/.style={circle, inner sep=0pt, minimum size=5pt, draw=none},
  figptoabierto/.style={circle, inner sep=0pt, minimum size=5pt, draw=figblue, line width=0.9pt, fill=white},
  figptocerrado/.style={circle, inner sep=0pt, minimum size=5pt, draw=figblue, line width=0.9pt, fill=figblue},
  figptoY/.style={figpto, fill=figblue},
  figptoZ/.style={figpto, fill=figred},
  % segmentos gruesos para intervalos sobre la recta
  figintervalo/.style={figblue, line width=2.2pt},
  % vectores/flechas de transformación (traslación, dilatación)
  figvec/.style={figred, line width=1pt, -{Latex[length=2.2mm,width=1.6mm]}},
  % etiquetas
  figlbl/.style={font=\small},
  figlblaux/.style={font=\footnotesize, figgray},
}

% ---- pgfplots: estilo base de las gráficas de funciones ----
\pgfplotsset{
  calcfig/.style={
    width=10cm, height=6cm,
    axis lines=middle,
    xlabel style={font=\small, at={(axis description cs:1,0.5)}, anchor=west},
    ylabel style={font=\small, at={(axis description cs:0.5,1)}, anchor=south},
    tick label style={font=\footnotesize},
    every axis plot/.append style={line width=1.1pt, figblue},
    grid=none,
    legend cell align=left,
    legend style={font=\footnotesize, draw=figgray!40},
    clip=false,
  },
}

% (compilar desde el directorio de la clase para que la ruta relativa resuelva).
%
% Paleta propia de este curso (no se reusa la de courses/Fisica3-2015): mismos
% tonos base (azul/rojo/ámbar) para alinear con keybox/notebox/warnbox, pero
% archivo independiente, sin dependencia cruzada entre cursos.
% ---------------------------------------------------------------------------

% ---- Paleta (alineada con las cajas del preámbulo: keybox/notebox/warnbox) ----
\definecolor{figblue}{HTML}{1F4E79}   % azul principal (líneas, keybox, conjuntos)
\definecolor{figred}{HTML}{B03A2E}    % rojo de énfasis (warnbox, región resaltada)
\definecolor{figamber}{HTML}{B8860B}  % ámbar (notebox)
\definecolor{figgray}{HTML}{555555}   % auxiliares, ejes, guías, universo

% ---- TikZ: librerías y estilos reutilizables ----
% Nota: las puntas se declaran como `-{Latex[...]}` y no como `->`. La forma
% con llaves NO contiene un `>` literal, así que es inmune al `>` activo que
% introduce babel español (de todos modos se carga \usetikzlibrary{babel} en
% el bloque §6.3 para cubrir cualquier `->` suelto).
\usetikzlibrary{arrows.meta,calc,angles,quotes,patterns}

\tikzset{
  figline/.style={figblue, line width=0.9pt},
  figaux/.style={figgray, dashed, line width=0.6pt},
  figaxis/.style={figgray, line width=0.7pt, -{Latex[length=2.2mm,width=1.5mm]}},
  % conjuntos (diagramas de Venn / patatoidales)
  figset/.style={figline, fill=figblue!6},
  figsetB/.style={figred, line width=0.9pt, fill=figred!6},
  figsetC/.style={figamber, line width=0.9pt, fill=figamber!10},
  figuniverso/.style={figgray, line width=0.7pt},
  % puntos de la recta / conjuntos numéricos
  figpto/.style={circle, inner sep=0pt, minimum size=5pt, draw=none},
  figptoabierto/.style={circle, inner sep=0pt, minimum size=5pt, draw=figblue, line width=0.9pt, fill=white},
  figptocerrado/.style={circle, inner sep=0pt, minimum size=5pt, draw=figblue, line width=0.9pt, fill=figblue},
  figptoY/.style={figpto, fill=figblue},
  figptoZ/.style={figpto, fill=figred},
  % segmentos gruesos para intervalos sobre la recta
  figintervalo/.style={figblue, line width=2.2pt},
  % vectores/flechas de transformación (traslación, dilatación)
  figvec/.style={figred, line width=1pt, -{Latex[length=2.2mm,width=1.6mm]}},
  % etiquetas
  figlbl/.style={font=\small},
  figlblaux/.style={font=\footnotesize, figgray},
}

% ---- pgfplots: estilo base de las gráficas de funciones ----
\pgfplotsset{
  calcfig/.style={
    width=10cm, height=6cm,
    axis lines=middle,
    xlabel style={font=\small, at={(axis description cs:1,0.5)}, anchor=west},
    ylabel style={font=\small, at={(axis description cs:0.5,1)}, anchor=south},
    tick label style={font=\footnotesize},
    every axis plot/.append style={line width=1.1pt, figblue},
    grid=none,
    legend cell align=left,
    legend style={font=\footnotesize, draw=figgray!40},
    clip=false,
  },
}


% ---- Hyperlinks & TOC ----
\usepackage[hidelinks]{hyperref}
\usepackage{bookmark}

% ---- Enumerate ----
\usepackage{enumitem}

% ---- Miscellaneous ----
\usepackage{icomma}
\usepackage{siunitx}
\usepackage{textcomp}
\usepackage[bottom]{footmisc}

% ---- Title ----
\title{\textbf{Clase 5: Orden en los reales, signos e intervalos} \\[0.3em]
        \large{Construcción de $\mathbb{N}$ por inducción}}
\author{Transcripto y expandido --- Curso Cálculo 1 (OpenFING)}
\date{}

\begin{document}

\maketitle
\thispagestyle{empty}
\newpage

\tableofcontents
\newpage

\section{Recapitulación: $\mathbb{R}$ es un cuerpo}

La clase retoma la lista de 15 axiomas de $\mathbb{R}$ presentada la clase
anterior, de la cual ya se habían visto 14. El \textbf{grupo 1} (axiomas 1.1
a 1.4) es el de la suma: \textbf{asociativa}, \textbf{conmutativa}, con
\textbf{neutro} $0$ y con \textbf{opuesto} $-x$ único para cada $x$. A partir
del opuesto se define la resta:
$$x - y := x + (-y).$$

El \textbf{grupo 2} (axiomas 2.1 a 2.5) es el del producto: asociativo,
conmutativo, \textbf{distributivo} respecto de la suma (la novedad frente al
grupo 1), con neutro $1 \neq 0$, y cada $x \neq 0$ tiene \textbf{inverso}
único $x^{-1}$. El cociente se define análogamente:
$$\frac{x}{y} := x \cdot y^{-1}.$$

Cuando un conjunto tiene una suma y un producto que cumplen los grupos 1 y 2,
se dice que es un \textbf{cuerpo}: se puede sumar, restar, multiplicar y
dividir sin restricciones. $\mathbb{R}$ es un cuerpo, pero no es el único
ejemplo. El docente repasa candidatos propuestos por la clase:

\begin{itemize}
  \item $\mathbb{Z}_2 = \{0,1\}$ con la suma módulo 2 ($0+0=0$, $0+1=1+0=1$,
    $1+1=0$) y el producto usual: es efectivamente un cuerpo (el más pequeño
    posible), pero no se trabajará con él porque le falta el orden total.
  \item $\mathbb{Q}$: cumple exactamente la misma lista de axiomas, luego es
    cuerpo.
  \item $\mathbb{C}$: también cumple la lista, luego es cuerpo.
\end{itemize}

\begin{notebox}[Cuerpos vistos en el curso]
En este curso solo se usarán $\mathbb{Q}$ y $\mathbb{R}$. Se menciona que
existen cuerpos finitos con 3, 5, 7, 11 elementos (en general, con $p$
elementos si $p$ es primo), pero no se estudiarán.
\end{notebox}

\section{Axiomas de orden total}

El \textbf{grupo 3} agrega una relación de orden $\leq$ (orden
\textbf{amplio}) sobre $\mathbb{R}$, con cinco axiomas:

\begin{enumerate}
  \item \textbf{Totalidad (tricotomía comparativa)}: para todo $x,y$, o bien
    $x \leq y$ o bien $y \leq x$ (siempre son comparables).
  \item \textbf{Transitividad}: $x \leq y$ y $y \leq z$ implican $x \leq z$,
    para todo $x,y,z$.
  \item \textbf{Antisimetría}: $x \leq y$ y $y \leq x$ implican $x = y$.
  \item \textbf{Compatibilidad con la suma}: $x \leq y$ implica
    $x + z \leq y + z$, para todo $x,y,z$.
  \item \textbf{Compatibilidad con el producto de no negativos}: si
    $x \geq 0$ y $y \geq 0$, entonces $xy \geq 0$.
\end{enumerate}

Estos cinco axiomas, junto con los 9 algebraicos de los grupos 1 y 2, forman
la lista completa de 14 axiomas vista hasta ahora. Falta un decimoquinto
(sección \ref{sec:completitud}), que es el que realmente distingue a
$\mathbb{R}$ de $\mathbb{Q}$.

\section{Orden amplio y orden estricto}

El orden \textbf{estricto} se define a partir del amplio:
$$x < y \; :\Longleftrightarrow\; x \leq y \ \text{y} \ x \neq y.$$

Recíprocamente, el orden amplio se puede recuperar del estricto:
$$x \leq y \; \Longleftrightarrow\; x < y \ \text{o} \ x = y.$$

\begin{notebox}[¿Cuál orden es el primitivo?]
Podría parecer más natural tomar $\leq$ como abreviatura de ``$<$ o $=$''
(el nombre lo sugiere), y de hecho se podría construir toda la teoría
axiomatizando el orden estricto en vez del amplio; se obtendría una
presentación completamente equivalente, solo cambiando la forma de los
axiomas. Lo importante es tener claras las propiedades de ambos, no cuál se
toma como primitivo.
\end{notebox}

\section{Reglas de las desigualdades}

\subsection{Suma: una traslación}

De los axiomas de orden y su combinación con la suma se obtiene una
\textbf{equivalencia lógica} (no solo una implicación en un sentido):
$$x \leq y \;\Longleftrightarrow\; x + z \leq y + z \;\Longleftrightarrow\; x - z \leq y - z, \qquad \forall x,y,z.$$

\textbf{Interpretación geométrica}: sumar $z$ a ambos miembros es
\textbf{trasladar} los dos puntos $x$ e $y$ de la recta por el mismo
desplazamiento $z$; como es el mismo movimiento para los dos, el orden
relativo entre ellos no cambia. Esto vale sea cual sea la posición de $0$
respecto de $x,y$.

También se tiene, tomando el opuesto:
$$x \leq y \;\Longleftrightarrow\; -y \leq -x.$$

\begin{warnbox}[El opuesto invierte el orden]
El paso al opuesto invierte el orden. Si $x$ está a la izquierda de $y$ en
la recta, $-x$ queda a la derecha de $-y$. Esta propiedad no depende de las
posiciones relativas de $x$, $y$ respecto de $0$. Las mismas dos
equivalencias valen reemplazando $\leq$ por $<$.
\end{warnbox}

La moraleja: sumar (o restar) una cantidad a una desigualdad nunca da
problemas: es una equivalencia lógica, y el razonamiento se puede invertir
en ambos sentidos.

\subsection{Producto: dilatación y contracción}

Multiplicar ambos miembros de una desigualdad por $z$ sí es delicado: el
resultado depende del \textbf{signo de $z$}.

\begin{itemize}
  \item Si $z > 0$: $x \leq y \implies xz \leq yz$ (no se invierte el
    orden).
  \item Si $z < 0$: $x \leq y \implies xz \geq yz$ (se invierte el orden).
  \item Si $z = 0$: $x \leq y \implies xz \leq yz$ es cierto pero inútil,
    porque siempre da $0 \leq 0$: se pierde información y no se puede
    recuperar la desigualdad original dividiendo, porque no existe
    $0^{-1}$.
\end{itemize}

$$
\boxed{
\begin{aligned}
z>0 &: \quad x \leq y \iff xz \leq yz \\
z<0 &: \quad x \leq y \iff xz \geq yz
\end{aligned}}
$$

La implicación recíproca (ir ``de vuelta'', dividiendo por $z$) solo es
lícita cuando $z$ es estrictamente positivo o estrictamente negativo, nunca
cuando $z=0$: en ese caso cualquier par $x \ne y$ colapsa a $0 \leq 0$, y no
hay forma de reconstruir la desigualdad de partida.

\textbf{Lectura geométrica} (para $z$ positivo): multiplicar por $z>1$ es
una \textbf{dilatación} del segmento $[x,y]$; multiplicar por $0<z\leq 1$
es una \textbf{contracción}. Para $z$ negativo se obtiene el mismo efecto
de escala pero además con \textbf{intercambio de posiciones} (lo que estaba
más a la izquierda pasa a estar más a la derecha):

\begin{table}[H]
\centering
\begin{tabular}{@{}llc@{}}
\toprule
Caso & Efecto & ¿Intercambia posiciones? \\
\midrule
$z \geq 1$ & dilatación & no \\
$0 \leq z \leq 1$ & contracción & no \\
$z \leq -1$ & dilatación & sí \\
$-1 \leq z \leq 0$ & contracción & sí \\
\bottomrule
\end{tabular}
\end{table}

% Clase 5 --- Interpretación geométrica de las reglas de las desigualdades
% (§4.1--4.2). El docente dibuja explícitamente estos tres casos en el
% pizarrón ("se puede interpretar en términos de traslación...", "tenemos
% dos dibujos posibles [dilatación/contracción]...", "voy a intercambiarse").
% Misma geometría base (recta con x<y) en los tres paneles; solo cambia la
% operación aplicada.
\begin{center}
\begin{tikzpicture}[scale=1]

  % ---------- Panel (a): Traslación (suma) ----------
  \begin{scope}[yshift=0cm]
    \draw[figaxis] (-0.8,0) -- (5.6,0);
    \node[figptocerrado] (x) at (0.5,0) {};
    \node[figptocerrado] (y) at (2,0) {};
    \node[figlbl, below=2pt] at (0.5,-0.05) {$x$};
    \node[figlbl, below=2pt] at (2,-0.05) {$y$};
    \node[figptocerrado] (xz) at (3,0) {};
    \node[figptocerrado] (yz) at (4.5,0) {};
    \node[figlbl, below=2pt] at (3,-0.05) {$x+z$};
    \node[figlbl, below=2pt] at (4.5,-0.05) {$y+z$};
    \draw[figvec] (x.north) to[bend left=35] node[figlbl, above, pos=0.5] {$+z$} (xz.north);
    \draw[figvec] (y.north) to[bend left=35] node[figlbl, above, pos=0.5] {$+z$} (yz.north);
    \node[figlbl, figblue, right, align=left] at (5.7,0)
      {(a) \textbf{traslación}: mismo\\ desplazamiento $z$ para\\ ambos, el orden se preserva};
  \end{scope}

  % ---------- Panel (b): Dilatación (producto, z>0) ----------
  \begin{scope}[yshift=-2.6cm]
    \draw[figaxis] (-0.8,0) -- (5.6,0);
    \node[figpto, fill=figgray] (o) at (0,0) {};
    \node[figlblaux, below=2pt] at (0,-0.05) {$0$};
    \node[figptocerrado] (x) at (1,0) {};
    \node[figptocerrado] (y) at (2,0) {};
    \node[figlbl, below=2pt] at (1,-0.05) {$x$};
    \node[figlbl, below=2pt] at (2,-0.05) {$y$};
    \node[figptocerrado] (xz) at (3,0) {};
    \node[figptocerrado] (yz) at (5,0) {};
    \node[figlbl, below=2pt] at (3,-0.05) {$xz$};
    \node[figlbl, below=2pt] at (5,-0.05) {$yz$};
    \draw[figvec] (x.north) to[bend left=20] node[figlbl, above, pos=0.5] {$\times z$} (xz.north);
    \draw[figvec] (y.north) to[bend left=20] node[figlbl, above, pos=0.5] {$\times z$} (yz.north);
    \node[figlbl, figblue, right, align=left] at (5.7,0)
      {(b) \textbf{dilatación} ($z>1$):\\ ambos se alejan del $0$,\\ el orden se preserva};
  \end{scope}

  % ---------- Panel (c): Contracción con inversión (producto, z<0) ----------
  \begin{scope}[yshift=-5.4cm]
    \draw[figaxis] (-3.2,0) -- (3.2,0);
    \node[figpto, fill=figgray] (o) at (0,0) {};
    \node[figlblaux, below=2pt] at (0,-0.05) {$0$};
    \node[figptocerrado] (x) at (1,0) {};
    \node[figptocerrado] (y) at (2,0) {};
    \node[figlbl, below=2pt] at (1,-0.05) {$x$};
    \node[figlbl, below=2pt] at (2,-0.05) {$y$};
    \node[figptoZ] (xz) at (-1,0) {};
    \node[figptoZ] (yz) at (-2,0) {};
    \node[figlbl, figred, below=2pt] at (-1,-0.05) {$xz$};
    \node[figlbl, figred, below=2pt] at (-2,-0.05) {$yz$};
    \draw[figvec] (x.north) to[bend left=30] (xz.north);
    \draw[figvec] (y.north) to[bend left=50] (yz.north);
    \node[figlbl, figred] at (0,1.55) {$\times z$};
    \node[figlbl, figred, right, align=left] at (3.3,0)
      {(c) $z<0$: los puntos\\ \textbf{cruzan} el $0$ y el\\ \textbf{orden se invierte}\\ ($xz>yz$)};
  \end{scope}

\end{tikzpicture}\\[0.5em]
{\small\itshape Misma recta base en los tres paneles ($x<y$): sumar $z$
(a) traslada ambos puntos por igual y nunca cambia el orden; multiplicar por
$z>0$ (b) los aleja o acerca del origen sin cambiar el orden; multiplicar por
$z<0$ (c) los hace cruzar el origen e \textbf{invierte} el orden --- el
origen de la mayoría de los errores con desigualdades.}
\end{center}


\begin{notebox}[Geometría antes que álgebra]
El docente insiste en que estas ideas geométricas (traslación para la suma,
dilatación/contracción con o sin intercambio para el producto) son más
importantes que la manipulación simbólica: ``no van a sobrevivir en
análisis si piensan que la matemática solo consiste en hacer cuentas sin
entender el contenido geométrico''. El curso podría llamarse igual de bien
``geometría analítica con números''.
\end{notebox}

\subsection{Cociente}

Las mismas reglas valen para el cociente entre ambos miembros de una
desigualdad (con $z \ne 0$):
$$
z>0: \; x \leq y \iff \frac{x}{z} \leq \frac{y}{z}, \qquad
z<0: \; x \leq y \iff \frac{x}{z} \geq \frac{y}{z}.
$$

\begin{notebox}[Ejercicio propuesto en clase]
Reescribir todas estas reglas con el orden estricto $<$. Cambian poco: los
casos con $z=0$ simplemente desaparecen (no tienen sentido para $<$, porque
$z=0$ fusiona a $x$ e $y$ en $0$), y en el resto basta reemplazar $\leq$ por
$<$ y $\geq$ por $>$.
\end{notebox}

\section{El inverso y el orden}

Invertir ($x \mapsto x^{-1}$) es, como el opuesto, una operación que puede
intercambiar el orden --- pero con una condición extra crucial: ambos
números deben ser no nulos y tener el mismo signo.

$$
\boxed{0 < x \leq y \;\Longrightarrow\; \frac{1}{y} \leq \frac{1}{x}}
\qquad\text{y análogamente para } x \leq y < 0.
$$

Cuando se cumple esta hipótesis (mismo signo, ambos no nulos), la relación
también se puede invertir de nuevo (aplicando el inverso una segunda vez se
recupera la desigualdad original).

\textbf{Verificación geométrica} con ejemplos concretos:

\begin{itemize}
  \item Positivos: $2 \leq 3 \implies \tfrac12 \geq \tfrac13$ (el orden se
    invierte).
  \item Negativos: $-3 \leq -2 \implies -\tfrac13 \geq -\tfrac12$ (también
    se invierte).
  \item Signos distintos: $-2 \leq 3 \implies -\tfrac12 \leq \tfrac13$ ---
    acá el orden se mantiene, porque el inverso de un negativo sigue siendo
    negativo y el de un positivo sigue siendo positivo, así que
    $1/x < 0 < 1/y$ automáticamente.
\end{itemize}

% Clase 5 --- El inverso y el orden (§5). El docente dibuja explícitamente
% "un dibujo entre dos casos" tomando 0,1,2,3 sobre la recta para comparar
% x, y y sus inversos, y remarca que hay que "saber posicionar
% intuitivamente" 1/x y 1/y. Esta figura junta ambos casos (mismo signo)
% en una sola recta con dos ramas.
\begin{center}
\begin{tikzpicture}[scale=1]
  % recta con las dos ramas
  \draw[figaxis] (-3.4,0) -- (3.4,0);
  \node[figpto, fill=figgray] (o) at (0,0) {};
  \node[figlblaux, below=4pt] at (0,-0.05) {$0$};
  \node[figlblaux] at (-2.3,0.55) {$\mathbb{R}^-$};
  \node[figlblaux] at (2.3,0.55) {$\mathbb{R}^+$};

  % rama positiva: x > 0 y su inverso 1/x, mismo signo
  \node[figptocerrado] (xp) at (2.4,0) {};
  \node[figptocerrado] (ixp) at (0.8,0) {};
  \node[figlbl, below=2pt] at (2.4,-0.05) {$x$};
  \node[figlbl, below=2pt] at (0.8,-0.05) {$\dfrac{1}{x}$};
  \draw[figvec] (xp.north) to[bend left=45] (ixp.north);
  \node[figlbl, figred, fill=white, inner sep=1.5pt] at (1.6,0.85) {inverso};

  % rama negativa: x < 0 y su inverso 1/x, mismo signo
  \node[figptocerrado] (xn) at (-2.4,0) {};
  \node[figptocerrado] (ixn) at (-0.8,0) {};
  \node[figlbl, below=2pt] at (-2.4,-0.05) {$x$};
  \node[figlbl, below=2pt] at (-0.8,-0.05) {$\dfrac{1}{x}$};
  \draw[figvec] (xn.north) to[bend right=45] (ixn.north);
  \node[figlbl, figred, fill=white, inner sep=1.5pt] at (-1.6,0.85) {inverso};

  \node[figlbl, figblue, align=center] at (0,-1.7)
    {si $x$ y $1/x$ tienen el \textbf{mismo signo}, ninguno cruza el $0$:\\
     $1/x$ queda en la \textbf{misma rama} que $x$ (más lejos si $|x|<1$, más cerca si $|x|>1$)};
\end{tikzpicture}\\[0.5em]
{\small\itshape En ambos casos (positivo y negativo) el inverso permanece
del mismo lado del $0$: por eso la comparación $0<x\leq y \implies 1/y \leq
1/x$ (y su análoga para negativos) tiene sentido sin cruces. Cuando $x,y$
tienen \textbf{signos distintos}, en cambio, el orden entre $1/x$ y $1/y$ se
\textbf{mantiene} automáticamente porque $1/x<0<1/y$.}
\end{center}


\begin{warnbox}[Cuidado con el inverso]
Cada vez que aparece un inverso en una desigualdad hay que verificar dos
cosas antes de concluir nada: que ambos lados son no nulos (si no, la
expresión ni siquiera está bien definida) y que tienen el mismo signo (si
no, no hay regla fija de inversión salvo el caso ``signos distintos''
descrito arriba). El docente lo remarca como fuente típica de errores:
``cuidado, hay un montón de trampas posibles, lo único posible es hacer
pequeños dibujos''.
\end{warnbox}

\section{Signos: $\mathbb{R}^+$, $\mathbb{R}^-$ y la regla de los signos}

Para todo $x \in \mathbb{R}$ vale la \textbf{tricotomía}: exactamente una de
las tres opciones se cumple,
$$x < 0 \quad \text{o bien} \quad x = 0 \quad \text{o bien} \quad x > 0.$$

Esto parte a la recta real en tres piezas: la parte negativa, el cero, y la
parte positiva. Se definen por comprensión:
$$\mathbb{R}^+ = \{x \in \mathbb{R} : x > 0\}, \qquad \mathbb{R}^- = \{x \in \mathbb{R} : x < 0\},$$
con la \textbf{descomposición fundamental}
$$\mathbb{R} = \mathbb{R}^- \cup \{0\} \cup \mathbb{R}^+.$$

\begin{notebox}[Convención sobre el signo de 0]
El docente cuenta que el cero tardó siglos en ser aceptado (``es muy suave
con la adición pero es la bomba atómica con la multiplicación''). Sobre la
convención de signo de $0$: en Francia se lo considera ni positivo ni
negativo (la descomposición de arriba, con $0$ excluido de ambos
conjuntos), mientras que en el mundo hispanohablante suele adoptarse la
convención inclusiva, donde ``positivo'' significa $x\geq 0$ y ``negativo''
significa $x \leq 0$, de modo que el $0$ tiene ``doble nacionalidad''. Es
solo una cuestión de convención, sin contenido matemático de fondo.
\end{notebox}

De los axiomas (no hay ningún axioma que lo postule directamente) se
\textbf{deducen} las reglas de los signos, tanto para la suma como para el
producto:

\begin{table}[H]
\centering
\begin{tabular}{@{}llll@{}}
\toprule
Suma & Resultado & Producto & Resultado \\
\midrule
$+ + +$ & $+$ & $+ \times +$ & $+$ \\
$+ + -$ & depende de cuál es mayor & $+ \times -$ & $-$ \\
$- + -$ & $-$ & $- \times +$ & $-$ \\
& & $- \times -$ & $+$ \\
\bottomrule
\end{tabular}
\end{table}

\begin{notebox}[Por qué se demuestra y no se memoriza]
Ninguno de los 14 axiomas menciona directamente ``negativo por negativo es
positivo''; es una consecuencia de las reglas algebraicas y de orden ya
vistas (ejercicio propuesto: demostrar que si $x<0$ e $y<0$ entonces
$xy>0$). El docente cuenta que este hecho --- contraintuitivo incluso para
niños --- fue descubierto por los matemáticos árabes hacia el año 1000
(aunque los números negativos ya existían antes, en India desde el siglo
VII y en China desde el siglo I, donde se entendían naturalmente en
términos de yin y yang), y que Europa tardó unos 700 años más, hasta el
siglo XVII, en aceptar los números negativos y el cero, por prejuicios
culturales (``obra del demonio'').
\end{notebox}

\section{Intervalos}

\subsection{Los diez tipos de intervalo}

Con el orden ya establecido se definen los intervalos por comprensión.
Dados $a \leq b$:
$$
[a,b] = \{x \in \mathbb{R} : a \leq x \leq b\}, \qquad
]a,b[ \ = \{x \in \mathbb{R} : a < x < b\},
$$
y de modo análogo los semiabiertos $[a,b[$ y $]a,b]$. Los semi-infinitos:
$$
\;]-\infty,a] = \{x : x \leq a\}, \quad ]-\infty,a[\ = \{x : x < a\},
$$
$$
[a,+\infty[\ = \{x : x \geq a\}, \quad ]a,+\infty[\ = \{x : x > a\}.
$$
A estos se agregan dos casos particulares: $]-\infty,+\infty[\ = \mathbb{R}$
(otro nombre, otra notación para $\mathbb{R}$) y el conjunto vacío
$\emptyset$, que también se considera un intervalo (aunque sin notación
propia como tal). En total, \textbf{diez tipos} de intervalo.

\begin{warnbox}[La hipótesis $a \leq b$ es esencial]
La definición de $[a,b]$ (o cualquiera de sus variantes) solo tiene el
sentido ``esperado'' cuando $a \leq b$. Si $a > b$, no existe ningún $x$
que sea a la vez $\geq a$ y $\leq b$, así que el conjunto que resulta de la
fórmula es simplemente $\emptyset$ --- vacío pero perfectamente bien
definido, no un error.
\end{warnbox}

% Clase 5 --- Los diez tipos de intervalo (§7.1). Catálogo visual de los
% cuatro casos más representativos: cerrado, abierto, semiabierto y no
% acotado. Misma geometría base (a=1, b=4) en las cuatro filas; solo cambia
% el tipo de extremo.
\begin{center}
\begin{tikzpicture}[scale=1]

  % ---------- Fila 1: cerrado [a,b] ----------
  \begin{scope}[yshift=0cm]
    \draw[figaux] (-0.6,0) -- (5.6,0);
    \draw[figintervalo] (1,0) -- (4,0);
    \node[figptocerrado] at (1,0) {};
    \node[figptocerrado] at (4,0) {};
    \node[figlbl, below=2pt] at (1,-0.05) {$a$};
    \node[figlbl, below=2pt] at (4,-0.05) {$b$};
    \node[figlbl, figblue, left] at (-0.8,0) {$[a,b]$};
  \end{scope}

  % ---------- Fila 2: abierto ]a,b[ ----------
  \begin{scope}[yshift=-1.6cm]
    \draw[figaux] (-0.6,0) -- (5.6,0);
    \draw[figintervalo] (1,0) -- (4,0);
    \node[figptoabierto] at (1,0) {};
    \node[figptoabierto] at (4,0) {};
    \node[figlbl, below=2pt] at (1,-0.05) {$a$};
    \node[figlbl, below=2pt] at (4,-0.05) {$b$};
    \node[figlbl, figblue, left] at (-0.8,0) {$\mathopen]a,b\mathclose[$};
  \end{scope}

  % ---------- Fila 3: semiabierto [a,b[ ----------
  \begin{scope}[yshift=-3.2cm]
    \draw[figaux] (-0.6,0) -- (5.6,0);
    \draw[figintervalo] (1,0) -- (4,0);
    \node[figptocerrado] at (1,0) {};
    \node[figptoabierto] at (4,0) {};
    \node[figlbl, below=2pt] at (1,-0.05) {$a$};
    \node[figlbl, below=2pt] at (4,-0.05) {$b$};
    \node[figlbl, figblue, left] at (-0.8,0) {$[a,b\mathclose[$};
  \end{scope}

  % ---------- Fila 4: no acotado [a,+infty[ ----------
  \begin{scope}[yshift=-4.8cm]
    \draw[figaux] (-0.6,0) -- (5.6,0);
    \draw[figintervalo, -{Latex[length=2.6mm,width=1.9mm]}] (1,0) -- (5.4,0);
    \node[figptocerrado] at (1,0) {};
    \node[figlbl, below=2pt] at (1,-0.05) {$a$};
    \node[figlbl, above=1pt] at (5.4,0.12) {$+\infty$};
    \node[figlbl, figblue, left] at (-0.8,0) {$[a,+\infty\mathclose[$};
  \end{scope}

\end{tikzpicture}\\[0.5em]
{\small\itshape Cuatro de los diez tipos de intervalo, misma geometría
($a=1$, $b=4$) en las cuatro filas: el círculo \textbf{lleno} marca un
extremo \textbf{incluido}, el círculo \textbf{vacío} uno \textbf{excluido};
la flecha abierta hacia $+\infty$ indica que no hay extremo superior real.
Los diez tipos surgen de combinar extremo izquierdo/derecho abierto o
cerrado, más los dos semi-infinitos, $\mathbb{R}$ y $\emptyset$.}
\end{center}


\subsection{Intersección de intervalos}

Se plantea calcular, en general, la intersección $[a,b] \cap [c,d]$ (con
$a\leq b$, $c \leq d$). El docente hace notar que la respuesta depende de
las posiciones relativas de $a,b,c,d$ (si $b < c$ da vacío; si se solapan,
da otra cosa; etc.), y da la fórmula única que cubre todos los casos:

$$
\boxed{[a,b] \cap [c,d] =
\begin{cases}
[\max(a,c), \min(b,d)] & \text{si } \max(a,c) \leq \min(b,d) \\
\emptyset & \text{en caso contrario}
\end{cases}}
$$

% Clase 5 --- Intersección de intervalos (§7.2). El docente comenta la
% aplicación a la detección de colisiones de rectángulos en videojuegos vía
% esta misma fórmula; acá se ilustra el caso con solape parcial que motiva
% max(a,c)/min(b,d).
\begin{center}
\begin{tikzpicture}[scale=1]

  % ---- guías verticales punteadas (dibujadas ANTES que las etiquetas,
  %      para que las etiquetas con fill=white las tapen limpiamente) ----
  \draw[figaux] (0,1.6) -- (0,-1.6);
  \draw[figaux] (5,1.6) -- (5,-1.6);
  \draw[figaux] (3,1.6) -- (3,-1.6);
  \draw[figaux] (8,1.6) -- (8,-1.6);

  % ---------- Fila superior: [a,b] ----------
  \begin{scope}[yshift=1.6cm]
    \draw[figaux] (-0.6,0) -- (8.6,0);
    \draw[figintervalo] (0,0) -- (5,0);
    \node[figptocerrado] at (0,0) {};
    \node[figptocerrado] at (5,0) {};
    \node[figlbl, below=2pt, fill=white, inner sep=1.5pt] at (0,-0.05) {$a$};
    \node[figlbl, below=2pt, fill=white, inner sep=1.5pt] at (5,-0.05) {$b$};
    \node[figlbl, figblue, left] at (-0.8,0) {$[a,b]$};
  \end{scope}

  % ---------- Fila media: [c,d] ----------
  \begin{scope}[yshift=0cm]
    \draw[figaux] (-0.6,0) -- (8.6,0);
    \draw[figintervalo] (3,0) -- (8,0);
    \node[figptocerrado] at (3,0) {};
    \node[figptocerrado] at (8,0) {};
    \node[figlbl, below=2pt, fill=white, inner sep=1.5pt] at (3,-0.05) {$c$};
    \node[figlbl, below=2pt, fill=white, inner sep=1.5pt] at (8,-0.05) {$d$};
    \node[figlbl, figblue, left] at (-0.8,0) {$[c,d]$};
  \end{scope}

  % ---------- Fila inferior: intersección ----------
  \begin{scope}[yshift=-1.6cm]
    \draw[figaux] (-0.6,0) -- (8.6,0);
    \draw[figred, line width=2.2pt] (3,0) -- (5,0);
    \node[figptoZ] (mc) at (3,0) {};
    \node[figptoZ] (mb) at (5,0) {};
    \node[figlbl, figred, below=2pt, anchor=north east, xshift=-3pt,
          fill=white, inner sep=1.5pt] at (3,-0.05) {$\max(a,c)$};
    \node[figlbl, figred, below=2pt, anchor=north west, xshift=3pt,
          fill=white, inner sep=1.5pt] at (5,-0.05) {$\min(b,d)$};
    \node[figlbl, figred, left] at (-0.8,0) {$[a,b]\cap[c,d]$};
  \end{scope}

\end{tikzpicture}\\[0.5em]
{\small\itshape Caso con solape parcial ($a<c<b<d$): el intervalo
intersección (fila roja) empieza en el \textbf{mayor} de los dos extremos
izquierdos, $\max(a,c)=c$, y termina en el \textbf{menor} de los dos
extremos derechos, $\min(b,d)=b$. Si $\max(a,c) > \min(b,d)$ (los
intervalos no se solapan) la intersección es directamente $\emptyset$.}
\end{center}


\begin{notebox}[Aplicación: intersección de rectángulos]
Ejercicio propuesto: demostrar la fórmula considerando todos los casos
posibles de posición relativa. El docente comenta su aplicación práctica en
informática: al programar videojuegos, decidir si dos rectángulos (p. ej.
dos personajes) se superponen se reduce a intersectar sus intervalos de
coordenadas $x$ e $y$ con exactamente esta fórmula. Advierte que ``la
intersección de intervalos es mucho más complicada de lo que uno podría
pensar'' --- hay muchos casos a distinguir, de los cuales en clase solo se
verificaron dos.
\end{notebox}

\section{Cuerpo totalmente ordenado y el axioma que falta}
\label{sec:completitud}

Cuando un conjunto cumple simultáneamente los axiomas de cuerpo (grupos 1 y
2) y los de orden total compatible (grupo 3), se dice que es un
\textbf{cuerpo totalmente ordenado}. Es una etiqueta que resume de una sola
vez los 14 axiomas ya vistos: en vez de enumerarlos uno por uno, un
matemático simplemente dice ``$\mathbb{R}$ es un cuerpo totalmente
ordenado''.

$\mathbb{Q}$ también es un cuerpo totalmente ordenado (cumple exactamente
la misma lista de 14 axiomas). $\mathbb{C}$, en cambio, no lo es, porque no
admite un orden total compatible con su estructura de cuerpo; tampoco lo es
$\mathbb{Z}_2$.

Esto deja una pregunta abierta: si $\mathbb{Q}$ y $\mathbb{R}$ satisfacen
exactamente la misma lista de 14 axiomas, ¿qué los distingue? La respuesta,
presentada como \emph{teaser} (adelanto sin desarrollar), es el
\textbf{decimoquinto axioma}, el \textbf{axioma de completitud}:

\begin{keybox}[Axioma de completitud (adelanto)]
Todo subconjunto no vacío de $\mathbb{R}$ y superiormente acotado tiene un
supremo.
\end{keybox}

$\mathbb{Q}$ no cumple este axioma (de ahí la diferencia esencial entre
ambos cuerpos), pero las nociones de ``superiormente acotado'' y
``supremo'' no se definen todavía --- quedan para la clase siguiente ---;
solo se identifican las palabras nuevas del enunciado (todo el resto,
``subconjunto'' y ``no vacío'', ya es conocido).

\section{Construcción de $\mathbb{N}$}

Con la estructura de cuerpo totalmente ordenado ya establecida, la clase
pasa a construir formalmente algunos subconjuntos notables de
$\mathbb{R}$, comenzando por $\mathbb{N}$.

\subsection{La función sucesor}

Se define la función \textbf{sucesor} $s: \mathbb{R} \to \mathbb{R}$,
$$s(x) := x + 1.$$

Intuitivamente, $\mathbb{N}$ se construye a partir de $0$ aplicando el
sucesor repetidas veces: $1 = s(0) = 0+1$, $2 = s(1) = 1+1$,
$3 = s(2) = 2+1$, etc. La definición ``de toda la vida''
$\mathbb{N} = \{0,1,2,\ldots\}$ es calificada de ``floja'': casi todos los
elementos quedan escondidos detrás de los puntos suspensivos. El objetivo
de la clase es dar una definición sin puntitos.

\subsection{Conjuntos inductivos}

\begin{keybox}[Definición: conjunto inductivo]
Un subconjunto $A \subseteq \mathbb{R}$ es \textbf{inductivo} cuando cumple
dos condiciones:
\begin{enumerate}
  \item $0 \in A$.
  \item $A$ es \textbf{estable por sucesor}: $x \in A \implies x+1 \in A$.
\end{enumerate}
\end{keybox}

Ejemplos discutidos en clase:

\begin{table}[H]
\centering
\begin{tabularx}{\textwidth}{@{}lcX@{}}
\toprule
Conjunto & ¿Inductivo? & Razón \\
\midrule
$[0,+\infty[$ & sí & contiene $0$; si $x\geq 0$ entonces $x+1\geq 0$ \\
$\{0,\tfrac12,1,\tfrac32,2,\ldots\}$ & sí & contiene $0$; sumar $1$ a cualquier elemento da otro elemento de la lista \\
$\mathbb{Z}$ & sí & contiene $0$ y es estable por sucesor \\
$\mathbb{Q}$ & sí & contiene $0$; la suma de un racional con $1$ es racional \\
positivos estrictos ($x>0$) & no & no contiene $0$ \\
$[0,42]$ & no & $42$ está, pero $s(42)=43$ no está \\
$[0,41]$ & no & $41$ está, pero $s(41)=42$ no está \\
\bottomrule
\end{tabularx}
\end{table}

\begin{notebox}[Cuidado circular]
El propio ejemplo de $\{0,\tfrac12,1,\tfrac32,\ldots\}$ generó una
discusión en clase: un estudiante lo propuso como conjunto inductivo, a lo
que el docente respondió que, efectivamente, lo es, pero que describirlo
así (con puntos suspensivos, generándolo paso a paso) ya usa implícitamente
la idea de inducción que se busca fundamentar, no presuponer.
\end{notebox}

\subsection{Definición de $\mathbb{N}$ sin puntos suspensivos}

Con la noción de conjunto inductivo ya disponible, se define:

$$
\boxed{n \text{ es un \textbf{entero natural}} \;:\Longleftrightarrow\; n \text{ pertenece a \emph{todos} los subconjuntos inductivos de } \mathbb{R}}
$$

y en consecuencia:
$$
\mathbb{N} := \{x \in \mathbb{R} : \forall A \subseteq \mathbb{R},\ (A \text{ inductivo} \implies x \in A)\}.
$$

Intuitivamente, $\mathbb{N}$ es la intersección de todos los subconjuntos
inductivos de $\mathbb{R}$: se toman todos los candidatos (el docente ya
exhibió varios) y se queda solo lo que está en común a todos ellos. Es una
definición más abstracta que la de ``puntitos'', pero formal y sin
ambigüedad.

\subsection{$\mathbb{N}$ es el subconjunto inductivo más pequeño}

\begin{keybox}[Proposición]
$\mathbb{N}$ es inductivo, y además es el más pequeño de los conjuntos
inductivos en el sentido de la inclusión: $\mathbb{N} \subseteq A$ para
todo $A$ inductivo.
\end{keybox}

\textbf{Demostración.}

\emph{Parte 1: $\mathbb{N}$ es inductivo.}
\begin{itemize}
  \item $0 \in \mathbb{N}$: por definición de conjunto inductivo, $0$
    pertenece a todo conjunto inductivo $A$; por lo tanto $0$ cumple la
    condición que define a $\mathbb{N}$ (pertenecer a todos los
    inductivos), luego $0 \in \mathbb{N}$.
  \item $\mathbb{N}$ estable por sucesor: sea $x \in \mathbb{N}$, es decir,
    $x$ pertenece a todo conjunto inductivo $A$. Como cada $A$ es, en
    particular, estable por sucesor, se tiene $x+1 \in A$ para ese mismo
    $A$. Como esto vale para cualquier inductivo $A$, se concluye que $x+1$
    pertenece a todos los inductivos, es decir, $x + 1 \in \mathbb{N}$.
\end{itemize}

\emph{Parte 2: $\mathbb{N} \subseteq A$ para todo $A$ inductivo.} Es
inmediata de la definición: $\mathbb{N}$ se definió exactamente como el
conjunto de los elementos que están en todos los inductivos, así que en
particular está contenido en cada uno de ellos. $\blacksquare$

\begin{notebox}[Frase para memorizar]
``$\mathbb{N}$ es el subconjunto inductivo más pequeño de $\mathbb{R}$''
resume, de una sola frase, tanto la definición formal como su relación con
todos los demás conjuntos inductivos vistos en los ejemplos.
\end{notebox}

\section{Principio de inducción completa}

\subsection{Enunciado}

Sea $P(n)$ un enunciado (una propiedad, a veces verdadera y a veces falsa)
sobre un natural cualquiera $n$. Si se cumplen:

\begin{enumerate}
  \item \textbf{Caso base}: $P(0)$ es verdadero.
  \item \textbf{Paso inductivo}: para todo $n \in \mathbb{N}$,
    $P(n) \implies P(n+1)$.
\end{enumerate}

entonces:
$$
\boxed{\forall n \in \mathbb{N},\ P(n) \text{ es verdadero.}}
$$

\subsection{Intuición: los dominós}

La imagen que propone el docente es la de una fila de dominós: si el
primero cae (caso base) y la distancia entre dominós consecutivos es tal
que la caída de uno siempre provoca la caída del siguiente (paso
inductivo), entonces todos los dominós de la fila terminan cayendo.

\begin{notebox}[Los dos ingredientes son necesarios]
Sin caso base ``nadie empuja el primer dominó y no cae ninguno''; y si en
algún punto se rompe la implicación (algún dominó no logra tumbar al
siguiente), la caída se detiene ahí y no se puede concluir nada sobre los
que siguen.
\end{notebox}

\subsection{Demostración}

A diferencia de la intuición informal, el principio de inducción se
demuestra a partir de la definición de $\mathbb{N}$ como el subconjunto
inductivo más chico (sección anterior); no es un axioma adicional.

\textbf{Demostración.} Se define
$$A := \{n \in \mathbb{N} : P(n) \text{ es verdadero}\}$$
(el ``conjunto de los dominós que caen''). Se muestra que $A$ es inductivo:

\begin{itemize}
  \item $0 \in A$: por hipótesis (1), $P(0)$ es verdadero, así que $0$
    cumple la condición que define a $A$.
  \item $A$ estable por sucesor: sea $x \in A$, es decir, $P(x)$ es
    verdadero. Por hipótesis (2), $P(x) \implies P(x+1)$, luego $P(x+1)$
    también es verdadero, es decir, $x + 1 \in A$.
\end{itemize}

Por lo tanto $A$ es un conjunto inductivo. Pero $\mathbb{N}$ es el
subconjunto inductivo más pequeño, así que $\mathbb{N} \subseteq A$. Por
otro lado, por construcción $A \subseteq \mathbb{N}$ (los elementos de $A$
son naturales que además cumplen $P$). De ambas inclusiones,
$A = \mathbb{N}$, es decir:
$$\forall n \in \mathbb{N},\ n \in A \iff \forall n \in \mathbb{N},\ P(n) \text{ verdadero.} \qquad \blacksquare$$

\begin{notebox}[No es magia]
El punto conceptual central: la inducción no cae del cielo ni es un axioma
separado; es una consecuencia directa de que $\mathbb{N}$ se definió como
el inductivo más pequeño.
\end{notebox}

\subsection{Variante desde un $n_0$ arbitrario}

Se enuncia (sin demostrar en clase, aunque se indica que la demostración es
análoga) una variante útil: sea $P(n)$ una propiedad y $n_0 \in \mathbb{N}$
fijo. Si $P(n_0)$ es verdadero y, para todo $n \geq n_0$,
$P(n) \implies P(n+1)$, entonces $P(n)$ es verdadero para todo $n \geq
n_0$.

\begin{notebox}[Misma idea, otro punto de partida]
Es exactamente el mismo principio, pero arrancando la cadena de dominós en
$n_0$ en vez de en $0$ --- útil porque en la práctica muchas inducciones
empiezan en $1$, $2$ o cualquier otro valor conveniente, no necesariamente
en $0$.
\end{notebox}

\bigskip
\noindent\emph{La próxima clase (lunes) retomará el axioma de completitud
--- definiendo ``superiormente acotado'' y ``supremo'' --- y se verán
ejemplos concretos de demostraciones por inducción.}

\end{document}
